\subsection{Critérios de Orientação Exigida}
\label{sec:criterios_orientacao_exigida}

Certas tarefas de \emph{rendezvous and docking/berthing}, não basta que o efetuador alcance a posição do alvo; é necessário que a sua orientação esteja alinhada dentro de uma tolerância determinada.  
Esse tipo de requisito acontece quando o engajamento exige que o efetuador esteja rotacionado de forma compatível com a interface mecânica de acoplamento.

De forma geral, esse critério pode ser expresso como:
\begin{equation}
\label{eq:workspace_orientado}
\angle(\mathbf{R}_E(q), \mathbf{R}_{req}) \leq \theta_{req},
\end{equation}

onde:
\begin{itemize}
    \item $\mathbf{R}_E(q)$: Matriz de rotação associada à configuração $q$ do efetuador;
    \item $\mathbf{R}_{req}$: Matriz de rotação de referência exigida pela tarefa;
    \item $\theta_{req}$: Limite máximo admissível para o erro angular.
\end{itemize}

Quando esse critério é imposto, define-se o chamado \emph{workspace orientado}, que é o subconjunto do workspace posicional $W_{pos}$ que satisfaz simultaneamente restrições de posição e de orientação.

\begin{equation}
W_{ori} \subseteq W_{pos}.
\end{equation}

No presente trabalho, workspace orientado não será considerado, pois o processo de captura não requer orientação específica do efetuador.
